\documentclass{umbruch-position}

\positionsnummer{ST-001 ES}
\title{Executive Summary}
\untertitel{System-theoretical analysis of recourse anxiety --- ST-001 v0.15}
\author{Um:bruch -- Denkfabrik f\"ur gesellschaftlichen Wandel}
\date{April 2026}
\versionsnummer{1.2}

\zusammenfassung{%
The economic efficiency review (\textgerman{Wirtschaftlichkeitspr\"ufung}, \S\S\,106\,ff. SGB\,V) system generates a recourse anxiety whose consequential costs could exceed the enforced recourse claims by a factor of 100 to 1,000, according to model calculations. The companion study ST-001 identifies 15 findings that clarify this discrepancy using historical, legal-dogmatic, behavioral-economic, and internationally comparative sources. \textbf{Genre and scope:} ST-001 is not an original study in the sense of new primary empirical data collection, but a \textbf{multi-perspective stock-taking with pilot character} --- a diagnostic construction (three certification errors) built on integrated legal-dogmatic, statistical, behavioral-economic, and institution-theoretical perspectives. 15 findings, three minimally invasive reform levers and 15 evaluated measures (cost--benefit table) offer concrete options for action for policymakers, associations, and civil society.%
}

\schlagworte{recourse, economic efficiency review, V2b, transparency, Executive Summary}

\begin{document}
\emergencystretch=3em
\maketitle

%% ============================================================
\section{Research Question}
%% ============================================================

The study investigates a striking discrepancy in the German economic efficiency review (\textgerman{Wirtschaftlichkeitspr\"ufung}, \S\S\,106\,ff. SGB\,V) system: The official recourse enforcement rate is below 1\,\% of statutory health insurance physicians --- yet 72\,\% of general practitioners report having experienced a recourse claim (\textgerman{Regress})\footnote{We use `recourse claim' for \textgerman{Regress} at first mention, subsequently abbreviated to `recourse'.} at least once (Ribbat et al. 2023, n=770), and 47\,\% state that the threat of recourse occupies their everyday practice ``strongly or very strongly''. The central question is: \textbf{What happens in the system that produces both simultaneously?} The study attempts to clarify this discrepancy using historical, legal-dogmatic, behavioral-economic, health-economic, and internationally comparative sources.

ST-001 is the scientific companion study to the position paper PP-003 (Recourse Transparency Portal).

%% ============================================================
\section{Key Findings}
%% ============================================================

The study identifies 15 essential findings, which are presented here with their respective evidence base and limitations.

% --------------------------------------------------
\subsection{Four measurement levels and their conflation}

\textbf{Finding:} The term ``recourse rate'' is used in public discussion for four fundamentally different metrics: audit rate, recourse proceeding rate, recourse enforcement rate, and recourse volume. These four levels are not comparable with each other, but are regularly conflated. The publicly communicated ``below 1\,\%'' figure refers to the recourse enforcement rate at the physician level, across both procedures (V1+V2), without procedural separation. The 99.7\,\% rejection rate (KV Hessen 2015) conversely measures the pipeline rejection rate of the anomaly audit (V1) alone.

\textbf{Evidence base:} KV statistics (KV Hessen 2015, Dt. \"Arzteblatt Art. 113630; KV BaW\"u 2022), proprietary systematic conceptual analysis.

\textbf{Limitation:} Nationwide, all four levels are missing as uniformly published figures. No Association of Statutory Health Insurance Physicians (KV) publishes statistics separated by procedure (V1 vs. V2).

% --------------------------------------------------
\subsection{Definitional divergence: Why ``below 1\,\%'' and ``72\,\%'' are not a contradiction (CORE4)}

\textbf{Finding:} The term ``recourse'' is defined differently by various actors. This definitional divergence explains the apparent discrepancy:

\smallskip
{\small
\begin{tabularx}{\textwidth}{l p{4.5cm} l l}
\hline
\textbf{Tier} & \textbf{What is counted} & \textbf{Number} & \textbf{Source} \\
\hline
1 (narrow) & Only V1, legally binding & ${<}$100/yr $\to$ 0.065\,\% & Virchowbund 2018 \\
2 (official) & V1+V2, legally binding & 21/21,500 $\to$ 0.097\,\% & KV BaW\"u 2022 \\
3 (de facto) & All V2 incl.\ settlements & approx.\ 47,000/yr $\to$ ${\sim}$30\,\% & BMG GVSG 2024 \\
\hline
\end{tabularx}
}
\smallskip

The factor between Tier~1 and Tier~3 is approx.\ 470--500$\times$. All actors are telling the truth --- about different metrics. The discrepancy is \textbf{not a perception error of the physicians, but a measurement problem of the statistics}. In Bavaria, the SpiFa reports 13,332 assessed recourses (2024) among $\sim$24,000 statutory health insurance physicians --- average amount 334\,EUR, $\sim$70\,\% of which are petty amounts under 300\,EUR.

\textbf{Evidence base:} BMG GVSG-Referentenentwurf 2024; SpiFa Bayern 2023/2024; Virchowbund 2018; KV BaW\"u 2022; multi-agent research (15 agents, 25 data points, 14 domains).

\textbf{Limitation:} The ``${\sim}$30\,\%'' estimate is based on the assumption of an even distribution of the 47,000 proceedings across 155,000 physicians. The actual distribution (by specialty, region, type of health insurance fund) is not published.

\textbf{Originality:} Following a systematic prior-art check (Opus agent, 20+ search queries), the synthesis --- definitional divergence as an explanation for the Ribbat discrepancy --- is \textbf{published nowhere}. All of the literature explains the discrepancy as a psychological problem of the physicians (availability bias, irrational fear). All individual building blocks (Ribbat, BMG, SpiFa, Virchowbund) are published; the synthesis is original.

% --------------------------------------------------
\subsection{Two audit domains (V1/V2) and their differing fairness}

\textbf{Finding:} The economic efficiency review system consists of two functionally incommensurable procedures:

\begin{itemize}
  \item \textbf{V1 (Anomaly audit, \S\,106 SGB\,V):} Statistical screening, initiated by the audit office, without individual suspicion. Broad filter with high sensitivity and low specificity --- a high rejection rate (99.7\,\%) is normal operation.
  \item \textbf{V2 (Individual case audit, \S\,106b SGB\,V):} Application-based suspicion procedure, initiated by the health insurance fund. Statutorily without ``advisory notice instead of recourse (\textgerman{Beratung statt Regress})'' (\S\,106b Abs.\,2 S.\,4 SGB\,V). The fund is simultaneously the applicant and the recipient of the recourse amount.
\end{itemize}

V1 has acquired protective tiers over the years (advisory notice before recourse 2012, practice-specific exemptions (\textgerman{Praxisbesonderheit}) 2011, target volumes abolished 2015). V2 explicitly remained without advisory protection --- this was not an oversight, but conscious legislative design.

\textbf{Evidence base:} Legislative texts (SGB\,V, GKV-VSG 2015), BSG case law, KBV statement on the GKV-VSG 2015, regional audit agreements.

\textbf{Limitation:} V2 data exists (BMG: approx.\ 47,000 proceedings nationwide in 2022; SpiFa Bavaria: 13,000 recourses/year), but is not published as an integrated enforcement statistic with rejection and settlement rates. No KV publishes pipeline data separated by procedure (V1 vs. V2).

% --------------------------------------------------
\subsection{V2a vs. V2b --- the crucial differentiation (CORE3), V2b internal differentiation (v0.14)}

\textbf{Finding:} The individual case audit (V2) breaks down into two subcategories with fundamentally different procedural logic, and V2b is further internally differentiated:

\begin{itemize}
  \item \textbf{V2a (Uneconomical prescription):} Guideline-compliant, but more expensive than necessary. Medical defense is possible; the cost difference method (\S\,106b Abs.\,2a SGB\,V, BSG 2024) limits the recourse to the difference to the economical alternative.
  \item \textbf{V2b (Inadmissible prescription):} Medical grounds are fundamentally irrelevant (normative concept of harm, BSG B\,6\,KA\,5/09\,R); no cost difference method; no curability after payment by the fund; no advisory protection. Internally differentiated into:
  \begin{itemize}
    \item \textbf{V2b\,(i) --- purely formal errors} (missing signature, stamp instead of signature, missing ICD, incorrect aut-idem cross): Formal conformity is not a proxy for medical quality. Full certification error\footnote{We use `certification error' for a structural flaw by which an official process (here: the economic efficiency review) claims to certify a quality it cannot validly measure.}.
    \item \textbf{V2b\,(ii) --- substantive inadmissibility} (Pharmaceutical Directive (\textgerman{Arzneimittel-Richtlinie}, AM-RL) of the Federal Joint Committee (G-BA) violation, off-label without justification, AMNOG exclusion): The formal rule already codes medical G-BA evidence; the certification error is weaker because the rule measures medical quality \emph{indirectly} --- the problem is the punitive communication instead of a learning signal.
  \end{itemize}
\end{itemize}

\textbf{Consequence for the overall picture:}
\begin{itemize}
  \item V1 is \textbf{fairer than portrayed in early versions of the study} (hearing, medical consideration, tier model).
  \item V2a has \textbf{defense options} (medical justification, cost difference).
  \item V2b is, \textbf{viewed in isolation, even more problematic} than previously portrayed: It is the only procedure in German law without the possibility of curing an error, without cost difference, and without medical relevance.
\end{itemize}

\textbf{Quantitative qualification (v1.2).} Not \emph{all} V2 proceedings are dysfunctional. An estimated \textbf{60--65\,\% of V2 proceedings} (Category~B, pure formal errors, plus procedure-external attack vectors such as the stamp recourse) fall under the second certification error; the remaining 35--40\,\% (Category~A: evidence-based formal rules --- AM-RL annexes, G-BA decisions, biosimilar requirement) encode medical evidence --- though via the punitive rather than the educational channel. The study's demand is therefore not the abolition of evidence-based steering, but the extension of the punitive channel with an educational one (advisory-notice protection also for V2, advance information, Lever~1).

\textbf{Evidence base:} GPE BaW\"u, KV Nordrhein, KV Bayern, BSG case law (B\,6\,KA\,5/09\,R, B\,6\,KA\,5/23\,R, B\,6\,KA\,27/12\,R), proprietary detailed research (CORE3 document).

\textbf{Limitation:} The V2a/V2b differentiation is not reflected in any official statistics. What proportion of V2 proceedings concerns V2a and what proportion concerns V2b is not publicly known.

% --------------------------------------------------
\subsection{Three certification errors --- twice in measurement content, once in measurement direction}

\textbf{Finding:} The recourse system exhibits three certification errors --- two in the measurement content, one overarching one in the measurement direction:

\begin{enumerate}
  \item \textbf{First certification error --- V1 (cost deviation $\neq$ medicine):} The anomaly audit measures cost deviation from the specialty group average, not medical quality. For a screening, this is technically correct; the error lies in the public interpretation of the low final rate as evidence of a ``functioning system''. The protective tier model filters almost entirely.
  \item \textbf{Second certification error --- V2b (form $\neq$ medicine):} The individual case audit claims to review medical uneconomical behavior. It de facto measures formal or regulatory vulnerability. The finding applies in full force to V2b\,(i) (purely formal errors); for V2b\,(ii) (substantive inadmissibility, AM-RL violation) it is mitigated because the formal rule encodes G-BA evidence --- the problem remains the punitive instead of educative communication.
  \item \textbf{Overarching certification error --- one-way audit (direction $\neq$ symmetry):} The entire system exclusively audits ``prescribed too much'', never ``under-supplied''. The care mandate of \S\,70 SGB\,V has no proprietary measuring instrument in the outpatient sector. The system thus causes a transfer of harm to the patient (medical + financial + temporal): When taking action, the physician is exposed to the risk of recourse; when the physician exercises restraint, the patient bears the under-supply, and in the case of evasive action by the physician, the over-supply (iatrogenic risks + financial/temporal copayments, travel distances, loss of working/vacation time). SVR trichotomy of \textgerman{Überversorgung, Unterversorgung, Fehlversorgung} (over-, under-, and misdirected care) (SVR 2000/2001, \S\,135a SGB\,V).
\end{enumerate}

Medical appropriateness within the meaning of \S\,12 Abs.\,1 SGB\,V and the care mandate of \S\,70 SGB\,V are systematically certified only outside the recourse system (arbitration boards of the medical associations, civil medical liability, no operative care audit in the outpatient sector).

\textbf{Evidence base:} Systematic analysis of all audit stages (pipeline diagrams); BSG case law; \S\S\,12, 70, 135a SGB\,V; Goodhart's Law / Campbell's Law; SVR expert report 2000/2001.

\textbf{Limitation:} The objection that the system has never made a medical quality promise is dogmatically correct (Objection 2 of the study). The certification error does not describe a flaw \textit{within} the system, but the gap between system measurement and public attribution of meaning. V1 considers medical specifics in the protective tier model; V2a allows medical defense; V2b\,(ii) codes G-BA evidence. The overarching directional error applies to \emph{all} procedures equally.

% --------------------------------------------------
\subsection{V1 protective tier model}

\textbf{Finding:} The anomaly audit (V1) contains a multi-stage protective tier model that has so far received little appreciation in the public debate:

\smallskip
{\small
\begin{tabularx}{\textwidth}{lX}
\hline
\textbf{Tier} & \textbf{Protective effect} \\
\hline
Hearing / statement & Physician can explain deviation (6-week deadline) \\
Practice-specific exemption (\textgerman{Praxisbesonderheit}) & Medical specifics of the patient structure are taken into account \\
Advisory notice instead of recourse (\textgerman{Beratung statt Regress}) & For first-time anomalies: advisory notice instead of recourse (\S\,106b Abs.\,2) \\
Waiting period (\textgerman{Karenz}) & Auditable only in the next period \\
Cap & 25,000\,EUR in the first two years after advisory notice \\
5-year rule & After 5 years without anomaly, one is considered a first-timer again \\
\hline
\end{tabularx}
}
\smallskip

The 99.7\,\% rejection rate is the result of a \textbf{functioning protective tier model}, not a dysfunctional filter.

\textbf{Evidence base:} GPE BaW\"u, KV Nordrhein, \S\,106b Abs.\,2 SGB\,V.

\textbf{Limitation:} All flagged physicians nevertheless bear the psychological costs of the screening (fear, effort), even the 99.7\,\% who are exonerated. Practice-specific exemptions (\textgerman{Praxisbesonderheit}) must be actively claimed --- those who do not know this lose.

% --------------------------------------------------
\subsection{Ribbat 72\,\% --- recourse experience without procedural separation}

\textbf{Finding:} 72\,\% of general practitioners and 78\,\% of orthopedists report having experienced a recourse at least once. 47\,\% (GP) / 55\,\% (Ortho) state that the threat of recourse occupies their everyday practice ``strongly or very strongly''. 51\,\% of general practitioners refer patients to specialists ``at least occasionally despite indication''. 85\,\% have withheld a prescription at least once out of fear of recourse.

\textbf{Evidence base:} Ribbat L, Linde K, Schneider A, Riedl B (2023): Gesundheitswesen 85(2): 111--118. DOI: 10.1055/a-1594-2527. Nationwide random sample, 1,000 general practitioners and 1,000 orthopedists each, response rate 41\,\% / 39\,\% (n=770). Chair of General Practice TU Munich.

\textbf{Limitation:} The study measures \textbf{V1+V2 together, without procedural separation}. It measures subjective perception and self-reported behavior, not objective care outcomes. An independent replication is pending. The causal conclusion ``audit pressure $\to$ under-supply'' is not empirically closed (Objection 1 of the study). The study is the best available data point, not conclusive proof.

\textbf{Computational plausibility check:} A Poisson model with three anchors ($\lambda = 0{,}30$ from the BMG federal calculation, $\lambda = 0{,}85$ from the KV Nordrhein payment rate, $\lambda = 1{,}48$ from the KV Nordrhein application rate) yields a bandwidth: A clear majority of physicians (60--99\,\% over three years, depending on the counted event type --- concluded proceeding, payment, or application) is affected. The 72\,\% Ribbat experience rate is thus computationally plausible, independent of the ``below 1\,\%'' enforcement rate (cf.\ ST-001 Sec.~3.4, summation calculation).

% --------------------------------------------------
\subsection{Game-theoretical Nash equilibrium}

\textbf{Finding:} Given a perceived recourse risk of approx. 15\,\% ($P_{\text{perceived}}$), ``not prescribing'' becomes the dominant strategy of the physician. The result is a Pareto-suboptimal Nash equilibrium: Physician loses (ethical stress, extra work), patient loses (under-supply), fund gains nothing (pays later for hospitalization), system loses the consequential costs. No one wants this equilibrium --- everyone plays it individually rationally.

Additionally, the study identifies a solidarity prisoner's dilemma within the medical profession: Whoever fights publicly becomes a target. Whoever remains silent benefits from the success of others.

\textbf{Evidence base:} Behavioral economic modeling based on the Ribbat data, Kahneman/Tversky loss aversion ($\lambda \approx 2{,}25$), KBV professional monitoring.

\textbf{Limitation:} The thresholds are model assumptions. The formal through-modeling with payoff matrices and sensitivity analyses is planned as a follow-up analysis, but is not yet available.

% --------------------------------------------------
\subsection{Consequential cost calculation (0.9--1.7 billion EUR, scenario assumption)}

\textbf{Finding:} The study estimates the consequential costs of recourse anxiety-induced misdirected care in the three-stage model:

\smallskip
{\small
\begin{tabularx}{\textwidth}{lXl}
\hline
\textbf{Stage} & \textbf{Value} & \textbf{Source} \\
\hline
System framework (ACSC + non-adherence, pre-attribution) & 7--17 bln. EUR/yr & WIdO, Zi, DIMDI HTA 65 \\
$\times$ Attribution rate (conservative) & 5\,\% & Scenario assumption \\
= Recourse consequential costs (overall bandwidth) & \textbf{0.9--1.7 bln. EUR/yr} & Proprietary projection \\
of which medium scenario & \textbf{approx.\ 1.1 bln. EUR/yr} & Sensitivity table ST-001 Sec.~4.3 \\
\hline
\end{tabularx}
}
\smallskip

This is contrasted by enforced recourses of an estimated 5--50 million EUR/year. \textbf{Scenario ratio of harm to revenue: approx.\ 100:1 to 1,000:1.}

In the medium scenario, approx.\ 6,230 preventable hospitalizations and approx.\ 125 deaths per year would be proportionally associable with recourse anxiety-induced under-supply.

\textbf{Evidence base:} WIdO Krankenhaus-Report, Zi 2021, DIMDI HTA 65, OECD Health at a Glance 2025, Ribbat 2023, proprietary projection with sensitivity table.

\textbf{Limitation:} \textbf{These numbers are an illustrative thought experiment based on scenario assumptions, not a causally-analytically secured finding.} In particular, the attribution rate (5\,\%) and the frequency of withheld prescriptions are model assumptions. None of the available studies directly measures these parameters. The incomplete denominator: The deterrent effect of the system (prevented additional expenditures) is real, but not quantified. The 100:1 figure depicts the worst case for the system, not the expected value.

% --------------------------------------------------
\subsection{Normative concept of harm and V2b incurability}

\textbf{Finding:} The normative concept of harm of the BSG (B\,6\,KA\,5/09\,R, 2010) defines: The harm to the fund is deemed to have occurred upon payment of the formally incorrect prescription --- regardless of whether the medication was medically indicated. Once proceedings have been initiated, the formal error can no longer be cured. The BSG confirmed this in 2025 (B\,6\,KA\,9/24\,R) with a 491,000\,EUR stamp recourse: The argument of medical indication was declared irrelevant.

In 2024, the BSG expressly excluded the difference cost regulation for ``inadmissible'' prescriptions (= formal errors) (B\,6\,KA\,5/23\,R; B\,6\,KA\,10/23\,R).

\textbf{Evidence base:} BSG case law (directly verified).

\textbf{Limitation:} In rare cases, courts have decided in favor of the physician despite formal errors (SG Mainz 2022: 260,000\,EUR rejected as an abuse of rights; LSG BaW\"u 2025: prohibition of double auditing). These exceptions prove by their rarity how little protection the legal route typically offers.

% --------------------------------------------------
\subsection{Misdirected care is bidirectional (under- AND over-supply) --- SVR trichotomy}

\textbf{Finding:} German recourse anxiety generates --- unlike US Tort Law --- misdirected care in two directions simultaneously (Terminology: SVR 2000/2001, \S\,135a SGB\,V):

\begin{itemize}
  \item \textbf{Under-supply in prescriptions} (\S\,70 SGB\,V): Physicians prescribe less, avoid expensive medications, forgo off-label therapies.
  \item \textbf{Over-supply in diagnostics and referrals} (\S\,12 SGB\,V): Physicians order additional blood tests, refer to specialists, request imaging --- not because the patient needs it, but for their own protection. (Goetz 2024: 54\,\% weekly unnecessary lab tests; Ribbat 2023: 51\,\% refer despite their own competence.)
\end{itemize}

The more precise term is \textbf{misdirected care} (\textgerman{Fehlversorgung}) as a cross-category: The system redirects medical resources from therapy to self-protection. This results in a \textbf{transfer of harm to the patient (medical + financial + temporal)} --- in both directions: (a) \textit{medically} iatrogenic risks and overdiagnosis burden; (b) \textit{financially and temporally} additional doctor appointments, copayments, waiting times, vacation days, travel distances, accompanying persons. What the audit system saves as ``economic efficiency'', the patient pays doubly --- as delayed therapy \emph{and} as the time/travel burden of unnecessary additional services.

\textbf{Evidence base:} Ribbat 2023, Goetz 2024, international comparison (Kessler \& McClellan 1996, Currie \& MacLeod 2008, Mello et al. 2010).

\textbf{Limitation:} Under-supply is methodologically much harder to measure than over-supply. When a physician withholds a necessary prescription, it does not appear anywhere --- except later in hospitalization statistics, without causal attribution.

% --------------------------------------------------
\subsection{V2 enforcement hidden in settlement (70--80\,\%)}

\textbf{Finding:} An estimated 70--80\,\% of physicians affected by V2 conclude settlements (typically 30--50\,\% of the recourse amount). The settlement is the enforcement --- only outside of court. It appears in the statistics as ``no recourse'', although the physician has paid. The rejection rate thus does not measure the quality of the applications, but the defensive strength of those affected.

Only two out of five termination paths produce a substantive finding (rejection notice or court judgment). Since almost all proceedings end before court, the system quality remains permanently unmeasured.

\textbf{Evidence base:} Practice reports from recourse consultants and medical professional associations, structural-procedural analysis.

\textbf{Limitation:} There is \textbf{no published study} on the settlement rate in V2. The 70--80\,\% estimate relies on practice reports, not systematic surveys.

% --------------------------------------------------
\subsection{Container genealogy (9 norm containers)}

\textbf{Finding:} In 1989, a single paragraph (\S\,106 SGB\,V) sufficed for the entire economic efficiency review. By 2026, this has mushroomed into five outpatient (\S\S\,106, 106a, 106b, 106c, 106d) and two inpatient (\S\S\,275, 275c) containers, \S\,48 BMV-\"A, plus framework guidelines and up to 17 regional audit agreements. The written regulatory surface area has roughly octupled in 35 years. The core regulatory substance has remained the same.

Three methodologically particularly revealing migrations: (1) The \S\,106a renaming trap (2017), (2) the cascading outsourcing of the individual case audit to three levels, (3) the parallel-opposing outpatient/inpatient operation. Seven of the ten identified structural inconsistencies in the legislative fabric could be fixed with purely redrafting effort. They have remained uncorrected for years.

\textbf{Evidence base:} Legislative texts (GRG 1989, GSG 1992, GKV-VSG 2015, TSVG 2019), BT-Drucksachen (Bundestag printed papers), historical versions via BGBl.

\textbf{Limitation:} The container analysis shows an increase in complexity, not a causal dysfunction. The practical consequences can only be derived via the findings from Part III (behavioral economics) and Part V (institutional silence) of the study.

% --------------------------------------------------
\subsection{International comparison}

\textbf{Finding:} Germany and Switzerland are practically the only countries in Western Europe with universal individual recourse-auditing targeted against individual physicians. France has been developing a structurally parallel debate under the term ``d\'elit statistique'' (MSO/MSAP procedures of the CNAM) since 2023. Other countries (UK, NL, Scandinavia) steer physicians' prescribing behavior via positive incentives (QOF), capitation, centralized pricing, or professional self-regulation --- without individual sanctions.

\textbf{Evidence base:} OECD Health at a Glance, NHS-QOF documentation, French primary sources (MG France, FMF, Le Quotidien du M\'edecin, CNAM control figures 2024).

\textbf{Limitation:} Comparability is limited by differing system architectures. The French case primarily concerns sick leaves, not pharmaceuticals.

% --------------------------------------------------
\subsection{Goodhart's Law / Campbell's Law}

\textbf{Finding:} The proxy failure of the audit system --- formal conformity is measured instead of medical quality --- is not an isolated case, but a known systemic pathology. Goodhart's Law (``When a measure becomes a target, it ceases to be a good measure'') and Campbell's Law describe exactly this mechanism. The audit rate becomes the target, formal conformity the object of measurement --- and both thereby lose their diagnostic power.

\textbf{Evidence base:} Theory-supported analysis, applied to the empirical findings from Parts I--III of the study.

\textbf{Limitation:} The application is analytical, not empirically tested.

% --------------------------------------------------
\subsection{Dual function: Evidence in formal disguise (CORE4)}

\textbf{Finding:} The analysis up to now treated ``formality'' as uniformly medically irrelevant. That was too blanket. The V2 audit triggers fall into three categories:

\smallskip
{\small
\begin{tabularx}{\textwidth}{l l X}
\hline
\textbf{Cat.} & \textbf{Share} & \textbf{Character} \\
\hline
A & ${\sim}$35\,\% & Rational evidence steering (biosimilar mandate, off-label, AM-RL) \\
B & ${\sim}$40\,\% & Purely formal control (stamp, ICD coding) \\
C & ${\sim}$25\,\% & Bureaucratic overhead (petty amounts ${<}$300\,EUR) \\
\hline
\end{tabularx}
}
\smallskip

The system has two parallel evidence channels: Continuing medical education obligations (CME, \S\,95d SGB\,V, \textit{preventative}) and individual case audit (V2, \textit{punitive}). The physician receives no learning signal prior to the recourse. The demand shifts: Not abolition of evidence steering (Cat.\,A), but supplementation of the punitive channel with an educative one.

\textbf{Evidence base:} AOK BaW\"u audit topic list 2025; SpiFa Bayern; AM-RL structure; \S\,95d SGB\,V.

\textbf{Limitation:} The percentages (35/40/25) are estimates based on the audit topic lists, not empirically measured. A breakdown separated by procedure into V2a/V2b does not exist.

\textbf{Originality:} Prior-art check (Opus agent, 20+ search queries): The synthesis --- recourse + CME as two parallel evidence channels (punitive vs. educative) --- is published nowhere. Stallberg, IQWiG, Zeschick, and Gandjour touch upon individual aspects; the synthesis is original.

% --------------------------------------------------
\subsection{Foucault's Panopticon}

\textbf{Finding:} The economic efficiency review system fulfills the three core features of Foucault's Panopticon: (1)~Asymmetrical visibility --- the physician does not know if and when their prescriptions will be audited, the fund sees their data in real-time. (2)~Internalization of control --- 72\,\% of general practitioners change their prescribing behavior preventatively, not because they are audited, but because they \textit{could} be audited. (3)~Economy of power --- the low V1 recourse rate is maximum efficiency: The screening alone generates widespread behavioral change without having to enforce many recourses.

The audit system is not merely a poor measurement instrument --- it is a highly effective disciplinary instrument that unfolds its greatest impact precisely \textit{because} it does not measure precisely.

\textbf{Evidence base:} Theory-supported analysis (Foucault, \textit{Discipline and Punish}, 1975), applied to the Ribbat findings and the structural-procedural analysis.

\textbf{Limitation:} Theoretical classification, not empirical proof of intentionality.

%% ============================================================
\section{Proposed Measures (CBR Table)}
%% ============================================================

The study evaluates 14 measures according to effort, political feasibility, expected benefit, and cost-benefit ratio (CBR). The ranking:

\smallskip
{\small
\begin{tabularx}{\textwidth}{c c X c c c c c}
\hline
\textbf{Rank} & \textbf{\#} & \textbf{Measure} & \textbf{Effort} & \textbf{Feasib.} & \textbf{Benefit} & \textbf{CBR} & \textbf{Time} \\
\hline
1 & 1 & \textbf{Half-sentence deletion \S\,106b Abs.\,2 S.\,4} (Advisory protection for V2) & low & high & high & \textbf{A+} & short \\
2 & 2 & Making \S\,48 BMV-\"A visible in the SGB\,V & low & high & med.\ & \textbf{A} & short \\
3 & 3 & Uniform nationwide petty limit & low & high & med.\ & \textbf{A} & short \\
4 & 4 & Positive feedback loop (clearance certificate) & low & high & med.\ & \textbf{A} & short \\
5 & 5 & Independent efficiency review of the audit system & low & high & high & \textbf{A} & med.\ \\
6 & 6 & Publication mandate for audit offices & low & med.\ & transf. & \textbf{A$-$} & short \\
7 & 7 & \textbf{PP-003: Recourse Transparency Portal} & med.\ & med.\ & high & \textbf{B+} & med.\ \\
8 & 8 & Opening of operative audit catalogs & med.\ & med.\ & high & B+ & med.\ \\
9 & 9 & Mandatory cost-consequence assessment for applications & med.\ & med.\ & med.\ & B & med.\ \\
10 & 10 & Maximum total proceeding duration & med.\ & med.\ & med.\ & B & med.\ \\
11 & 11 & Open-source prescription tool & med.\ & med.\ & med.\ & B & med.\ \\
12 & 12 & Sanction consequence for unfounded fund applications & med.\ & low & high & B & long \\
13 & 13 & FOIA queries to audit offices & low & low & low--m. & C & short \\
14 & 14 & Governance pharmacovigilance & high & low & transf. & B$-$ & long \\
15 & 15 & \textbf{Care audit (\S\S\,70, 135a SGB\,V)} --- measuring instrument for under-/misdirected care (symmetrical counter-audit to the economic efficiency review) & high & med.\ & transf. & B & long \\
\hline
\end{tabularx}
}
\smallskip

\textbf{Explanation of the most important measures:}

\begin{itemize}
  \item \textbf{Rank 1 (Half-sentence deletion):} One line of legislative text would eliminate the structural asymmetry between V1 and V2. This would mean that, as with the anomaly audit, an advisory notice instead of a recourse occurs on the first occasion for individual case audits as well. \textbf{Dogmatically important:} The advisory notice according to \S\,106b Abs.\,2 S.\,2 SGB\,V is a \textit{regulatory decision with the quality of an administrative act} (\textit{Regelungsentscheidung mit Verwaltungsaktqualit\"at}) (not just an informal hint; BSG B\,6\,KA\,21/15\,R, B\,6\,KA\,1/17\,R). It establishes grounds for legitimate expectation (\textit{Vertrauenstatbestand}) and precludes later recourses for the same facts. Already proposed in the KBV critique in 2015.
  \item \textbf{Rank 2 (\S\,48 BMV-\"A):} The most important existence-destroying reason for recourse (approx.\ 490,000\,EUR stamp recourse, BSG B\,6\,KA\,9/24\,R of 27.08.2025) lives outside the SGB\,V in a collective agreement. Whoever ``only'' reads the law overlooks it. Purely redrafting amendment, no actor loses.
  \item \textbf{Rank 5 (Efficiency review):} No actor (BRH, BAS, SVR, Bundestag) audits the economic efficiency of the economic efficiency review. A statutory evaluation mandate every 5 years would break the institutional dynamic of self-preservation.
  \item \textbf{Rank 6 (Publication mandate):} Without published enforcement statistics, any efficiency discussion is impossible. The game-theoretical analysis identifies the dismantling of asymmetric information as the most impactful single intervention.
\end{itemize}

%% ============================================================
\section{Significance for the Recourse Transparency Portal (PP-003)}
%% ============================================================

PP-003 positions itself in the upper midfield of the measure ranking (\textbf{Rank 7 of 14}, CBR: B+).

\textbf{Strategic function:} PP-003 is the ``Plan B intended to force Plan A.'' If the publication mandate of the audit offices (Rank 6) is realized politically, the portal will become partially redundant --- and exactly that would be the desired outcome. The portal generates the informational pressure; the reform levers generate the institutional pressure. One without the other fizzles out.

\textbf{What PP-003 can achieve:}
\begin{itemize}
  \item Provide procedurally separated (V1 vs. V2) data from the physician's perspective for the first time
  \item Correct the risk calibration of physicians (aggregate numbers instead of anecdotes)
  \item Build political reform pressure by making the data gap visible
  \item Build civil society pressure for the measures in Ranks 1--6
\end{itemize}

\textbf{What PP-003 cannot achieve:}
\begin{itemize}
  \item Eliminate the structural asymmetry between V1 and V2 (requires Rank 1)
  \item Substitute for the lack of efficiency oversight (requires Rank 5)
  \item Prevent the compliance theater on the side of the health insurance funds (Transparency Paradox according to Bernstein 2012)
  \item Empirically close the causal inference audit pressure $\to$ under-supply
\end{itemize}

\textbf{Robustness against critique:} PP-003 takes the Transparency Paradox (Bernstein 2012) into account through three design decisions: (1)~Voting booth principle (cryptographic separation), (2)~$k$-anonymity ($k\geq5$), (3)~Focus on aggregates instead of individual cases. Nevertheless, the portal is a necessary but not sufficient step.

%% ============================================================
\section{What should be implemented first}
%% ============================================================

\subsection*{Highest Priority (Ranks 1--6): Low effort, high feasibility, implementable in the short term}

Ranks 1--6 share three characteristics: They cost little, are politically connectable, and can become effective within one year.

\textbf{The most efficient single lever} is the half-sentence deletion (Rank 1): One line of legislative text would eliminate the single most important driver of recourse anxiety --- the fact that, unlike the anomaly audit, no ``advisory notice instead of recourse (\textgerman{Beratung statt Regress})'' is provided for the individual case audit.

\textbf{The most transformative lever} is the publication mandate for the audit offices (Rank 6): As long as no enforcement statistics exist, any reform debate is a debate without a data foundation.

\textbf{The politically simplest levers} are Ranks 2--4: purely redrafting amendments (making \S\,48 BMV-\"A visible), measures already included in KBV lists of demands (petty limit), or KV-internal measures without statutory changes (clearance certificate).

\subsection*{Medium Priority (Ranks 7--11): Supplementary measures with medium effort}

PP-003 (Rank 7) unveils its full impact only in combination with Ranks 1--6. It is particularly valuable as a civil society pressure build-up when the political levers are not pulled promptly.

\subsection*{Long-term Vision (Ranks 12--14): Structural reforms}

Governance pharmacovigilance (Rank 14) --- a regular, independent evaluation of the care-side effects of regulatory measures --- would be the most comprehensive system change. It has the highest transformative benefit, but the lowest political feasibility.

%% ============================================================
\section{Unresolved Questions}
%% ============================================================

\subsection*{6.1 Missing V2 pipeline data}

For the individual case audit (V2), no published pipeline statistic exists. Neither the total number of applications, nor the rejection rate, nor the settlement rate, nor the duration of proceedings is publicly published. The entire column ``n.\,D.'' (no data) in the pipeline tables of the study is waiting to be filled.

\subsection*{6.2 FOIA queries (Deadline approx.\ 2026-05-10)}

Nine Freedom of Information Act (FOIA/IFG) queries to BMG, BAS, KBV, GPE, and other offices were submitted on 10 April 2026. The response period is 30 days (\S\,7 Abs.\,5 IFG). Initial acknowledgments of receipt are available. The most important expected insights:

\begin{itemize}
  \item \textbf{A1 + A2 (BMG, GPE RLP):} Would fill the V2 pipeline with real numbers for the first time --- the ``largest single leap in quality'' for the study.
  \item \textbf{A3 (BAS):} Would empirically underpin or refute the Niskanen finding (GPE self-preservation).
  \item \textbf{A6 (BRH):} The response ``no special audit since 2005'' would be a finding in itself.
  \item \textbf{A9 (BMG):} Would force the BMG to provide a written statement on concrete weak points in the legislative fabric.
\end{itemize}

Refusals are also journalistically exploitable: ``This data is not available to us'' confirms the diagnosed transparency gap.

\subsection*{6.3 Causality gap (Audit pressure $\to$ under-supply)}

The causal conclusion ``audit pressure $\to$ changed prescribing behavior $\to$ under-supply $\to$ patient harm'' is not empirically closed. The study lacks a causal-analytical study that directly correlates care outcomes with audit intensity --- such as a comparison of KV regions with differing individual case audit frequencies. Such a natural experiment would be the strongest piece of evidence, but does not exist.

\subsection*{6.4 Ribbat replication pending}

Ribbat et al. (2023) is the only large German survey study that specifically measures the effects of recourse anxiety on prescribing behavior. An independent replication --- ideally with a procedurally separated survey (V1 vs. V2) --- is pending. Two empirical findings would significantly weaken the causal chain:

\begin{enumerate}
  \item A replication with less than 20\,\% reported behavioral effect
  \item A correlation analysis between audit intensity and prescription volume that shows no connection
\end{enumerate}

\subsection*{6.5 Consequential cost calculation not empirically founded}

The projection (0.9--1.7 billion EUR/year) is based on scenario assumptions. In particular, the attribution rate (5\,\%) and the frequency of withheld prescriptions are model assumptions. The model calculation does not replace a causal-analytical study. A systematic review by a responsible body (BRH, SVR, IGES) has not yet taken place.

\subsection*{6.6 Deterrent effect not quantified}

The study calculates only the costs (under-supply, defensive medicine, consequential harm), not the benefits of the audit system (prevented additional expenditure through deterrence). This deterrent effect is real and difficult to measure. A complete cost-benefit analysis would have to account for it in the denominator.

\subsection*{6.7 V2a/V2b distribution unknown}

What proportion of individual case audits concerns V2a (uneconomical) and what proportion concerns V2b (inadmissible/formal error) is not publicly known. Without this differentiation, the severity of the finding cannot be calibrated.

\subsection*{6.8 Automated screening by funds --- open legal question}

Whether health insurance funds actually systematically scan all prescription data with audit software to generate individual case audit applications from it is not empirically proven. It is proven that the technical means exist and that the volume (KV Nordrhein: approx.\ 20,000 applications/year) is difficult to explain solely by random suspicion. The constitutional classification (BVerfG Hessendata 2023, 1\,BvR\,1547/19) is argumentatively defensible, but not judicially clarified.

\subsection*{6.9 GPE costs externally unverifiable}

The costs of the nationwide audit infrastructure (approx.\ 600--900 positions) are not verifiable from open sources. The GPE are not state authorities, do not appear in any budget, and the Federal Court of Audit (BRH) does not audit them. This structural non-verifiability is a finding in itself.

%% ============================================================
\section{Methodological Transparency}
%% ============================================================

The study was created in a multi-stage, AI-supported analysis process in April 2026 (estimated total working time approx.\ 60--80 hours, human + AI combined). 31 independent research threads were conducted and transferred into an integrated argumentation.

\textbf{Multi-model architecture:} Four instances in division of labor --- Claude (Opus 4.6, primary analysis and synthesis), Copilot (GPT-4o, concept sharpening and misconception-detection), Gemini (Deep Research, theoretical framework and source supplementation), LG/Human (question formulation, directional decisions, final release).

\textbf{Bias management:} Cross-model divergences were utilized systematically as signals. Five documented hypothesis revisions (i.a., dating of the individual case audit: GSG 1992 $\to$ GRG 1988; Moosazadeh misattribution $\to$ Zheng 2023; cost model: system framework $\to$ three-stage model after Copilot review). A dedicated Devil's Advocate agent systematically attempted to refute the central thesis.

\textbf{Cross-source divergence (Phase XIV):} 15 specialized agents (6 Opus, 9 Sonnet) researched in parallel with 25 data points from 14 independent domains. Central lesson: Cross-model divergence controls model bias, but not source bias. Only multi-source research brought the SpiFa Bavaria and BMG data to light. Two prior-art checks confirmed: Both central syntheses (definitional divergence; evidence in formal disguise) are original.

\textbf{Hallucination controls:} Gemini outputs required systematic double-checks on current political events (documented cases: Moosazadeh misattribution, French protest numbers). Every source was checked against the original publication.

\textbf{Limits of evidence:} The study does \textbf{not} claim that the economic efficiency review system should be abolished, that all recourses are unjustified, or that physicians do not require auditing. It claims:

\begin{enumerate}
  \item The system predominantly does not measure what it claims to measure (certification error) --- whereby ${\sim}$35\,\% of V2 audits absolutely provide evidence-based steering. The problem is not the evidence, but the communication channel: punitive instead of educative.
  \item The individual case audit is structurally asymmetric in design (no advisory protection, normative concept of harm for V2b).
  \item The ``below 1\,\%'' figure and the 72\,\% experience rate measure different things (definitional divergence, factor ${\sim}$500$\times$). The discrepancy is a measurement problem, not a perception error.
  \item The data situation is so thin that neither advocates nor critics can empirically fully prove their positions.
  \item The transparency gap is not accidental, but structural --- and closing it would be the single most effective step toward objectifying the debate.
\end{enumerate}

Where the study makes scenario assumptions, they are marked as such. Where appraisals are made, they are recognizable as appraisals. Where findings allow multiple readings, all readings are presented.

\vspace{4mm}
\noindent\rule{\textwidth}{0.4pt}

\vspace{2mm}
\noindent\textbf{Complete study:} ST-001 System-theoretical analysis of recourse anxiety (Working version v0.14, April 2026)\\
\textbf{Companion study to:} PP-003 Recourse Transparency Portal (v2.9)\\
\textbf{License:} CC BY 4.0\\
\textbf{Contact:} \href{mailto:hallo@um-bruch.org}{hallo@um-bruch.org} \quad|\quad um-bruch.org\\
\textbf{Repository (sources, research, versions):} \href{https://github.com/research-line/regressangst}{github.com/research-line/regressangst}\\
\textbf{Companion software (Open Source):} \href{https://github.com/research-line/verordnungsampel}{github.com/research-line/verordnungsampel} \textit{-- VerordnungsAmpel, tool for recourse risk indicator display without warranty.}

\vspace{4mm}
\noindent\rule{\textwidth}{0.4pt}

\vspace{2mm}
\noindent\textbf{Change history:}\\
\textbf{v1.0} (April 2026) --- First version, based on ST-001 v0.11. 15 findings, CBR table (14 measures), PP-003 alignment, 9 open research questions.\\
\textbf{v1.0+CORE4} (April 2026) --- CORE4 integration: 2 new findings (definitional divergence with three-tier table; dual function with A/B/C categories). V2 data section corrected (BMG 47,000, SpiFa Bavaria 13,332). ``Limits of evidence'' expanded to 5 points. Multi-agent experiment + prior-art checks in Methodological Transparency. ST-001 reference updated to v0.12 (87~pages). 13~pages.\\
\textbf{v1.1+CORE4} (April 2026) --- Synchronization with ST-001 v0.14 (95~pages) after 18-fold review chain: Poisson bandwidth $\lambda = 0{,}30$--$1{,}48$ with 60--99\,\% overall bandwidth added; advisory notice as a regulatory decision (\S\,106b Abs.\,2 S.\,2 SGB\,V; BSG B\,6\,KA\,21/15\,R, B\,6\,KA\,1/17\,R) clarified; overarching certification error and SVR trichotomy (\S\S\,70, 135a SGB\,V) firmly anchored; care audit as new CBR row~15; consequential cost detailing to medium scenario $\approx 1{,}1$\,bln.\,EUR; BSG B\,6\,KA\,9/24\,R date (27.08.2025) and sum (approx.\ 490,000\,EUR) unified.\\
\textbf{v1.2} (April 2026, 2026-04-14) --- Synchronization with ST-001 v0.15 (P1 block): \textbf{Genre clarification} in the abstract --- ``multi-perspective stock-taking with pilot character'' instead of an original study; \textbf{quantitative qualification} in the V2a/V2b section: 60--65\,\% of V2 proceedings (Category~B + procedure-external) fall under the second certification error, 35--40\,\% (Category~A) encode G-BA evidence punitively rather than educationally (Devil's Advocate finding from Objection~5 promoted into the main thesis). Subtitle and ST-001 reference updated to v0.15.

\vspace{2mm}
\noindent\textbf{AI transparency notice:} This document was developed using extensive deployment of AI language models (i.a., Anthropic Claude, Microsoft Copilot, Google Gemini). Several control and review cycles were conducted~--- AI-supported and human. AI models are prone to hallucinations; errors cannot be excluded despite careful checking. We correct documented errors and publish improved versions. This document does not constitute legal advice. Feedback: \href{mailto:hallo@um-bruch.org}{hallo@um-bruch.org}

\end{document}
